\begin{figure}[!hbp]
\centering
\begin{tikzpicture}[x=11mm,y=9mm,font=\small,
  branch/.style={draw=ink!80,line width=.65pt},
  marked/.style={draw=sea,line width=1pt},
  cluster/.style={draw=ink!55,rounded corners=2pt,fill=white},
  meets/.style={draw=sea,rounded corners=2pt,fill=sea!12},
  block/.style={|-|,draw=ink,line width=.7pt}]
\foreach \i in {1,...,8}
  \coordinate (leaf\i) at (\i,0);
\coordinate (root) at (4.1,3.2);
\coordinate (left) at (2,2.05);
\coordinate (right) at (5.7,2.3);
\coordinate (leftpair) at (2.5,.85);
\coordinate (s) at (6.5,1.65);
\coordinate (pairfive) at (5.5,.75);
\coordinate (pairseven) at (7.5,.75);
\draw[branch] (left)--(root)--(right)
              (leaf1)--(left)--(leftpair)
              (leaf2)--(leftpair)--(leaf3)
              (leaf4)--(right)--(s);
\draw[marked] (pairfive)--(s)--(pairseven)
              (leaf5)--(pairfive)--(leaf6)
              (leaf7)--(pairseven)--(leaf8);
\foreach \n in {root,left,right,leftpair} \fill[ink] (\n) circle (1.4pt);
\foreach \n in {s,pairfive,pairseven} \fill[sea] (\n) circle (1.6pt);
\node[right=3pt,text=sea] at (s) {$s$};
\foreach \i in {1,...,8} {
  \fill[ink] (leaf\i) circle (1.4pt);
  \node[below=1pt] at (leaf\i) {\(\i\)};
  \draw[densely dotted,ink!45] (\i,-.52)--(\i,-1.92);
}
\foreach \k/\y in {2/-.9,4/-1.72}
  \node[anchor=east] at (.25,\y) {\(k=\k\)};
\draw[cluster] (.55,-.65) rectangle (3.45,-1.15);
\foreach \l/\r in {4/6,7/8}
  \draw[meets] (\l-.45,-.65) rectangle (\r+.45,-1.15);
\draw[cluster] (.55,-1.47) rectangle (4.45,-1.97);
\draw[meets] (4.55,-1.47) rectangle (8.45,-1.97);
\foreach \l/\r in {1/1,2/3,4/4,5/6,7/8}
  \draw[block] (\l-.30,-.9)--(\r+.30,-.9);
\foreach \l/\r in {1/3,4/4,5/8}
  \draw[block] (\l-.30,-1.72)--(\r+.30,-1.72);
\end{tikzpicture}
\caption{An ordered tree of formula splits, with slots numbered in leaf
order. The bars show the maximal subtrees of at most \(k\) leaves, and
the boxes group consecutive bars into clusters. The highlighted subtree
has \(B_s=\{5,6,7,8\}\). It meets two clusters when \(k=2\) and lies in
one when \(k=4\), as the shaded boxes show. All four leaves therefore
receive the same displacement at the larger scale.}
\label{fig:cluster}
\end{figure}
